\documentclass[11pt]{article}
\usepackage[a4paper,margin=1in]{geometry}
\usepackage{amsmath,amssymb}
\usepackage[hidelinks]{hyperref}
\pdftrailerid{}

\newcommand{\Q}{\mathbb Q}

\title{Certified Engine-A Addendum for RQ-000013\\
\large to the supplementary material for\\
\emph{Effective Archimedean Stark Theorems over Real Quadratic Fields}}
\author{Hainan Zhao}
\date{31 July 2026}

\begin{document}
\maketitle

\section*{Claim boundary}

The identity below is \texttt{PROVED}: the uniform Engine-A theorem is
applied after every case-specific hypothesis is checked by exact
arithmetic.  The accompanying \texttt{bnrL1} point-value comparison is
only an \texttt{OBSERVED} cross-check and is not used in the proof.
This versioned addendum is part of the correction release at
\url{https://doi.org/10.5281/zenodo.21712478}; it does not alter the
already-published Zenodo v1.3 files.

\section*{The first nonzero imprimitive branch}

Let \(K=\Q(\sqrt2)\), with finite ideal in the frozen integral basis
\[
 \mathfrak f=\begin{pmatrix}14&6\\0&2\end{pmatrix},
 \qquad N\mathfrak f=28,
\]
and modulus \(\mathfrak m=\mathfrak f\infty_2\).  Its ray group is
\(G\simeq C_2\), in pinned PARI coordinates \(G=\{[0],[1]\}\);
the sign class is \(R=[1]\).  The unique supported character has
\(\chi([1])=-1\).  Its primitive conductor is
\[
 \mathfrak f_\chi=\begin{pmatrix}7&3\\0&1\end{pmatrix}\infty_2.
\]
The unique omitted prime has norm \(2\), valuation \(2\) in
\(\mathfrak f\), valuation \(0\) in \(\mathfrak f_\chi\), and
\(\chi^\circ(\mathfrak p)=-1\).  Hence the exact imprimitive factor is
\[
 E_\chi=1-\chi^\circ(\mathfrak p)=2.
\]

The quadratic ray field is represented relatively and absolutely by
\[
 t^2-2\sqrt2+1,\qquad P(t)=t^4+2t^2-7.
\]
It has signature \((2,1)\), discriminant \(-448\), class number \(1\),
and two roots of unity; \texttt{bnfcertify} returns \(1\).
On its free-unit basis the exact norm map is
\[
 \begin{pmatrix}1&-1\end{pmatrix},
\]
whose primitive kernel is generated by \((1,1)^T\).  The embedded base
unit has coordinates \((1,-1)^T\), so
\[
 I_\chi=\left|\det
 \begin{pmatrix}1&1\\-1&1\end{pmatrix}\right|=2.
\]

Put
\[
 u=\frac{t^2-2t+3}{4}.
\]
The chosen absolute generator is the unique root of \(P\) in
\((-34/25,-27/20)\).  Exact Sturm counts and rational endpoint
evaluation give
\[
 \frac95<u<\frac{19}{10}.
\]
Thus this is the oriented relative unit.  It has minimal polynomial
\[
 U^4-2U^3+U^2-2U+1,
\]
and the nontrivial Artin element \(t\mapsto-t\) sends \(u\) exactly to
\(u^{-1}\).

Since \(h_K=h_L=1\), \(w_K=w_L=2\), \(E_\chi=2\), and \(I_\chi=2\),
the uniform theorem gives
\[
 L'_{\mathfrak m}(0,\chi)=2\log u.
\]
Here \(|G|=2\), so the Fourier coefficient \(2/|G|\) is one.  Therefore
the fully Artin-labelled packet is
\[
 \boxed{X_{[0]}=u^2,\qquad X_{[1]}=u^{-2}.}
\]
The selected value \(u^2\in(7/2,18/5)\) has minimal polynomial
\[
 X^4-2X^3-5X^2-2X+1.
\]

\section*{Replay record}

The source of truth is the JSON record whose repository-relative path
is split here only for typesetting:
\begin{center}
\ttfamily\small
artifacts/rq000013-engine-a-\linebreak
imprimitive-certificate-v1.json
\end{center}
Its deterministic replay command is
\begin{verbatim}
python3 scripts/certify_rq000013_engine_a.py --check
\end{verbatim}
under PARI/GP 2.15.4.  The frozen transcript is at
\begin{center}
\ttfamily\small
artifacts/rq000013-engine-a-\linebreak
imprimitive-certificate-v1.transcript.
\end{center}

\end{document}
